\documentclass[12pt,a4paper]{article}

\usepackage[utf8]{inputenc}
\usepackage[T5]{fontenc}
\usepackage[vietnamese]{babel}
\usepackage{geometry}
\usepackage{graphicx}
\usepackage{array}
\usepackage{booktabs}
\usepackage{longtable}
\usepackage{enumitem}
\usepackage{setspace}
\usepackage{hyperref}
\usepackage{xcolor}

\geometry{
	top=2.5cm,
	bottom=2.5cm,
	left=3cm,
	right=2cm
}

\onehalfspacing

\begin{document}
	
	% =========================
	% COVER PAGE
	% =========================
	
	\begin{titlepage}
		\begin{center}
			
			{\Large \textbf{University of Transport Ho Chi Minh City}}\\[0.5cm]
			
			{\large \textbf{FACULTY OF INFORMATION TECHNOLOGY}}\\[2cm]
			
			{\Large \textbf{CAPSTONE PROJECT}}\\[1cm]
			
			{\LARGE \textbf{
					AI-POWERED PERSONALIZED LEARNING PATH PLATFORM
					FOR COMPETENCY ASSESSMENT EXAM PREPARATION
			}}\\[2cm]
			
			\begin{flushleft}
				
				\textbf{Project Title:}
				AI-powered Personalized Learning Path Platform for
				Competency Assessment Exam Preparation
				
				\vspace{0.5cm}
				
				\subsection{Group Members}
				
				\begin{enumerate}
					\item Nguyễn Trung Nguyên -- 066205001323
					\item Trương Nguyễn Trí Cường -- 070206007667
					\item Nguyễn Thúy Nhi -- 082307011843
					\item Nguyễn Phạm Quỳnh Như -- 068306003038
					\item Phan Đình Quân -- 045206001313
				\end{enumerate}
				
			\end{flushleft}
			
			\vfill
			
			Ho Chi Minh City, 2026
			
		\end{center}
	\end{titlepage}
	
	% =========================
	% TABLE OF CONTENTS
	% =========================
	
	\tableofcontents
	
	\newpage
	
	% =========================
	% MAIN CONTENT
	% =========================
	
	\section{Main Proposal Content}
	
	\subsection{Context}
	
	In recent years, Competency Assessment Examinations have become
	one of the primary admission methods adopted by many universities
	in Vietnam.
	
	Unlike traditional examinations that mainly evaluate memorized
	knowledge, competency assessment focuses on measuring students'
	logical reasoning, analytical thinking, reading comprehension,
	quantitative reasoning, and problem-solving abilities.
	
	Although numerous online learning platforms and examination
	preparation systems are available, most of them provide identical
	learning content and predefined study plans for every learner.
	
	Such an approach does not consider individual differences in
	academic background, learning pace, target scores, available study
	time, or knowledge mastery.
	
	With the rapid advancement of Artificial Intelligence and
	Educational Data Analytics, personalized learning has become a
	promising approach for improving learning efficiency.
	
	Therefore, this project proposes the development of an AI-powered
	personalized learning platform that supports competency assessment
	exam preparation through intelligent study planning, adaptive
	practice recommendation, learning analytics, and AI-assisted
	tutoring.
	
	% =========================
	% PROPOSED SOLUTIONS
	% =========================
	
	\subsection{Proposed Solutions}
	
	The proposed system aims to build an intelligent learning platform
	that continuously analyzes students' academic performance and
	automatically personalizes their learning journey throughout the
	examination preparation process.
	
	The system provides the following capabilities:
	
	\begin{itemize}
		
		\item Managing learning materials, question banks, mock
		examinations, and competency frameworks.
		
		\item Conducting diagnostic assessments to identify each
		student's strengths and weaknesses.
		
		\item Generating personalized learning paths according to
		target scores and available preparation time.
		
		\item Recommending adaptive learning content and practice
		questions.
		
		\item Monitoring learning progress using learning analytics
		dashboards.
		
		\item Providing AI-powered tutoring for question explanation,
		concept clarification, and knowledge summarization.
		
		\item Predicting examination performance and recommending
		learning adjustments.
		
		\item Supporting teachers and parents in monitoring students'
		learning progress through analytical reports.
		
	\end{itemize}
	
	% =========================
	% SYSTEM ROLES
	% =========================
	
	\subsubsection{System Roles}
	
	\paragraph{Administrator}
	
	\begin{itemize}
		\item Manage users and system roles.
		\item Manage courses and competency frameworks.
		\item Manage question banks and learning resources.
		\item Configure AI services and system settings.
		\item Monitor overall platform operations.
	\end{itemize}
	
	\paragraph{Teacher / Academic Manager}
	
	\begin{itemize}
		\item Develop and organize learning materials.
		\item Manage knowledge domains and learning objectives.
		\item Create and maintain mock examinations.
		\item Evaluate question quality.
		\item Review students' learning progress.
		\item Monitor AI recommendation quality.
	\end{itemize}
	
	\paragraph{Student}
	
	\begin{itemize}
		\item Complete diagnostic assessments.
		\item Set target examination scores.
		\item Follow personalized learning paths.
		\item Practice adaptive exercises and mock examinations.
		\item View learning dashboards.
		\item Interact with AI Tutor.
		\item Receive personalized study recommendations.
		\item Track predicted examination performance.
	\end{itemize}
	
	\paragraph{Parent}
	
	\begin{itemize}
		\item Monitor students' learning progress.
		\item View competency development reports.
		\item Receive learning alerts and notifications.
		\item Track target score achievement.
		\item Monitor study schedules and examination readiness.
	\end{itemize}
	
	\paragraph{AI Recommendation Engine and Chat Bot}
	
	\begin{itemize}
		\item Analyze learning profiles.
		\item Predict examination performance.
		\item Explain examination questions.
		\item Provide detailed answer explanations.
		\item Answer students' academic questions.
		\item Generate additional practice questions.
	\end{itemize}
	
	% =========================
	% RBL
	% =========================
	
	\subsubsection{Research-Based Learning Topics}
	
	\begin{itemize}
		\item Adaptive Learning Systems
		\item Intelligent Tutoring Systems
		\item Educational Recommendation Systems
		\item Knowledge Tracing
		\item Student Performance Prediction
		\item Learning Analytics
		\item Educational Data Mining
		\item Large Language Models for Education
		\item Retrieval-Augmented Generation
	\end{itemize}
	
	% =========================
	% FUNCTIONAL REQUIREMENTS
	% =========================
	
	\subsection{Functional Requirements}
	
	\subsubsection{Core Flow 1: Competency Diagnosis and Learning Profile Generation}
	
	Students begin by taking a diagnostic assessment designed to
	evaluate their competency across multiple knowledge domains and
	cognitive skills.
	
	The system analyzes assessment results together with behavioral
	learning indicators to establish a comprehensive Learning Profile.
	
	\subsubsection{Core Flow 2: Personalized Learning Path Planning}
	
	Based on the generated Learning Profile, the student's target
	examination score, and the available preparation period, the AI
	engine constructs a personalized study roadmap.
	
	\subsubsection{Core Flow 3: Adaptive Learning and Continuous Recommendation}
	
	Throughout the learning process, students complete lessons,
	quizzes, and mock examinations recommended by the system.
	
	\subsubsection{Core Flow 4: AI-assisted Learning Support}
	
	Students can interact with an AI Tutor at any stage of their
	learning journey.
	
	\subsubsection{Core Flow 5: Learning Analytics and Progress Monitoring}
	
	The platform continuously collects learning activities and
	visualizes these indicators through interactive dashboards.
	
	\subsubsection{Core Flow 6: Examination Performance Prediction}
	
	The AI engine predicts each student's examination performance
	based on historical learning data, competency progression,
	and behavioral analytics.
	
	% =========================
	% NON FUNCTIONAL
	% =========================
	
	\subsection{Non-Functional Requirements}
	
	\begin{enumerate}
		
		\item Average response time below 3 seconds for standard operations.
		
		\item Support for concurrent access by at least 1,000 active users.
		
		\item Secure authentication using JWT with Role-Based Access
		Control (RBAC).
		
		\item Modular and scalable architecture separating Mobile,
		Web, Backend, AI Services, and Analytics components.
		
		\item Secure storage of learning data with scheduled backup
		and recovery mechanisms.
		
		\item Cloud-native deployment supporting high availability
		and horizontal scalability.
		
		\item APIs designed for future integration with Learning
		Management Systems (LMS) and additional examination frameworks.
		
	\end{enumerate}
	
	% =========================
	% TECHNOLOGY
	% =========================
	
	\subsection{Theory and Practical Technologies}
	
	\begin{table}[htbp]
		\centering
		\begin{tabular}{|p{5cm}|p{8cm}|}
			\hline
			\textbf{Component} & \textbf{Technology} \\
			\hline
			
			Mobile Application & Flutter (Android and iOS) \\
			\hline
			
			Web Application & ReactJS / Next.js \\
			\hline
			
			Backend Services & ASP.NET Core Web API / Node.js Express \\
			\hline
			
			Database & PostgreSQL \\
			\hline
			
			Authentication & JWT and OAuth2 \\
			\hline
			
			Large Language Models & LLMs \\
			\hline
			
			AI Integration & OpenAI GPT API \\
			\hline
			
			RAG & LangChain / Semantic Kernel \\
			\hline
			
			Vector Database & Qdrant / Pinecone \\
			\hline
			
			Cloud Deployment & Microsoft Azure \\
			\hline
			
			CI/CD & GitHub Actions \\
			\hline
			
			Version Control & GitHub \\
			\hline
			
		\end{tabular}
		
		\caption{Technology Stack}
		\label{tab:technology}
	\end{table}
	
	% =========================
	% PRODUCTS
	% =========================
	
	\subsection{Expected Deliverables}
	
	\begin{enumerate}
		
		\item Mobile application for students.
		
		\item Web portal for teachers and academic administrators.
		
		\item Parent portal for learning progress monitoring.
		
		\item Backend platform for learning data management.
		
		\item AI Recommendation Engine.
		
		\item AI Tutor integrated with LLM and RAG.
		
		\item Learning Analytics Dashboard.
		
		\item Examination Performance Prediction Module.
		
		\item Personalized Learning Path Management Module.
		
	\end{enumerate}
	
	% =========================
	% TASKS
	% =========================
	
	\subsection{Proposed Tasks}
	
	\begin{longtable}{|c|p{9cm}|c|}
		\hline
		\textbf{No.} & \textbf{Task} & \textbf{Member} \\
		\hline
		
		1 &
		Perform requirement analysis, system architecture design,
		database design, and backend API development
		& All \\
		\hline
		
		2 &
		Develop the Flutter mobile application for students,
		including personalized learning, adaptive practice,
		AI Tutor interaction, and learning progress tracking
		& 1 \\
		\hline
		
		3 &
		Develop the web portal for administrators, teachers,
		and parents, including learning content management,
		question bank administration, dashboards, reports,
		and user management
		& 2 \\
		\hline
		
		4 &
		Design and implement AI components, including the
		Recommendation Engine, Knowledge Tracing,
		Student Performance Prediction, and LLM-based
		AI Tutor with Retrieval-Augmented Generation
		& 3 \\
		\hline
		
		5 &
		Conduct system testing, performance optimization,
		cloud deployment, technical documentation,
		and final Capstone report preparation
		& All \\
		\hline
		
	\end{longtable}
	
	% =========================
	% OTHER COMMENTS
	% =========================
	
	\subsection{Other Comments}
	
	Additional comments and related proposals will be developed
	during the implementation and evaluation stages of the project.
	
	% =========================
	% END
	% =========================
	
\end{document}